% R124B prospective proof chain; do not relabel as part of F124A.
\section{Paid-value action generation and staged policy certificates}
\label{supp:valuebridge124b}

\begin{proof}[Proof of Proposition~\ref{prop:valuebridge124b}]
The defining strong-convexity inequality gives, for $u\in[0,1]$,
\[
 q((1-u)s+ut)
 \le(1-u)q(s)+uq(t)-\frac\mu2u(1-u)(t-s)^2.
\]
The two weights are nonnegative, so their endpoint upper enclosures
can be substituted to obtain \eqref{eq:valuebridge124b}.
For the stated information-class sharpness only, the polynomial
\[
 q_*(v)=U_s+\frac{U_t-U_s}{d}(v-s)
                   +\frac\mu2(v-s)(v-t)
\]
has second derivative $\mu$, takes endpoint values $U_s,U_t$,
and is admissible on $[0,A]$ for every $L\ge\mu$ when the two
level intervals and curvature are the only restrictions. Its
value at $b=s+ud$ is exactly $C_{st}(b)$. This proves sharpness
in that class without claiming compatibility with additional
measured derivatives or raw global supports.

As a function of $u$, the proposed upper bound equals
\[
 U_s+(U_t-U_s-\mu d^2/2)u+\mu d^2u^2/2.
\]
Its second derivative is $\mu d^2>0$. Its unconstrained
stationary point is $1/2-(U_t-U_s)/(\mu d^2)$, and projection
onto $[0,1]$ gives the unique constrained minimizer. No division
by zero occurs because $s<t$. The midpoint is feasible directly.
All operations are rational for rational inputs.

For a realized $b$ the proof uses its actual barycentric coefficient
$u=(b-s)/d$. If $b\notin[s,t]$, one coefficient is negative and
the substitution of both upper endpoint values is not justified;
this is why such a pair is excluded. If $s,t$ are grid points
and $s\le b^*\le t$, downward rounding on that same grid obeys
$s\le\lfloor b^*/h\rfloor h\le t$. General feasible rounding
needs only to stay in $[0,A]$; the upper evaluation then uses
the pairs that actually bracket the rounded action and the
always-valid fused bound. Thus the theorem neither assigns an
ideal proposal's value to a different action nor charges an
unproved derivative of the implemented future error.
\end{proof}

\paragraph{Validity of the shared original-q information.}
The paid risk supports are global original-q minorants.
The fused derivative envelopes enclose $q'$ throughout the
feasible interval, so their directed integrals give valid
original-q contrasts. Consequently a representative construction is
\begin{align*}
 \underline q(b)&=\max\left\{\max_k S_k(b),
             \max_i[W_i+C_F^-(a_i,b)]\right\},\\
 U_F(b)&=\min_i[U_i+C_F(a_i,b)],
\end{align*}
where $S_k$ are the saved global supports and $C_F^-$ is
the fused lower contrast. Every maximum or minimum is over
quantities with the same stated orientation. The common optimum
lower can be the maximum of the inherited qualified lower
and the all-action minimum of the maximum of valid supports.
Taking a restricted candidate minimum would not provide this
all-action lower bound. These constructions are the retained
fused service, with all original computation costs included.

Each bridge piece is an original-q upper on its own bracket,
so $U_B(b)$ is also a valid upper. Since
$q(a)\ge\underline q(a)$,
\[
 q(b)-q(a)\le U_B(b)-\underline q(a).
\]
Its minimum with the fused upper remains a valid contrast.
Also $q(b)\le q(a)+C_B(a,b)\le U_B(a)+C_B(a,b)$,
so the transported intersection $U_B^q(a,b)$ is a valid
absolute upper for the bridge stage.
The actual original node vector belongs to the joint program:
the inserted origin and candidate nodes carry unknown variables,
not invented observations. Its paid inequalities, enclosing
boxes and adjacent interpolation conditions hold for that
same original q. The rational dual theorem therefore gives
$q(b)-q(a_B)\le D(a_B,b)$. Intersecting with the cheap
contrast proves $C_J$. Finally
$q(b)\le q(a_B)+C_J(a_B,b)\le U_B(a_B)+C_J(a_B,b)$,
which proves the second upper intersection in
\eqref{eq:jointbridgecertificate124b}. Neither step requires
a fresh q value at the selected bridge action.

\begin{proof}[Proof of Theorem~\ref{thm:stagedpolicy124b}]
Keep the future controller fixed. For distinct feasible actions its
error satisfies $\eta(b)-\eta(a)\le\delta$.
Adding it to a valid q-contrast upper proves the first line
of \eqref{eq:stagedqualification124b}. The bounds
$0\le\eta(b)\le\delta$ and
$\underline q(b)\le q(b)\le U^q(b)$ prove the absolute
policy interval. Because $B\le V_0=\min q$, subtracting B
from the policy upper proves the regret upper. The nonnegative
regret follows from feasibility and the original nested recursion.
The two qualification predicates are separate: neither the
strict contrast sign nor the original-optimum gap implies a
narrow absolute interval on its own.

The fused service supplies a qualified starting action.
At the bridge stage, either that action is retained with
zero gain, or the three tests prove a feasible strictly
improved action with the same regret and width requirements.
At the joint stage exactly the same argument applies with
the bridge action as the incumbent. Its same-ledger contrast
and absolute bounds were established above. Thus induction
gives \eqref{eq:stagedorder124b} and both qualifications.
At identical actions use the policy identity instead of
paying a comparison allowance. Summing the two accepted
inequalities telescopes the actual policy values and proves
the cumulative gain statement.

The conclusion allows any finite fixed menu and deterministic
tie rule. The particular pair and neighbor menu is generated
before either enhancement is evaluated, ensuring that both
stages receive the same feasible options. The joint stage
can decline a candidate already selected by the bridge,
because retaining its incumbent is an identity. A sound
but loose numerical dual cannot cause a false gain; its
old bound is retained by the intersection. A recoverable
enhancement failure must retain the current qualified
incumbent, so the joint stage retains the bridge policy.
Source, model, state or infrastructure failures remain
failed executions under the protocol; they cannot silently
change the policy definition or sampling census.

Feasible rounding is part of the returned action. Every
q bound, risk support, pair membership, interpolation node,
and contrast in this proof is evaluated at that realized
action. A finite rational quantity can then be converted
outward to reported endpoints. Accepting only after the
outward regret and width tests preserves the inequalities;
no differentiability of a rounded action is required.
\end{proof}

\paragraph{Why these are complete policies rather than changed targets.}
The original one-date continuation $V_1=T_1g$ is acquired by paid
original-risk calculations and shared with the conventional controls.
The terminal controller and its uniform error are the same at every
candidate. None of the bridge algebra changes its law, feasible set,
discount or objective. The new operation changes the current feasible
action only after an original complete-policy comparison. A smaller
dual residual without an accepted action is a certificate improvement,
not an observed economic improvement. Identical-feature hinges have
the same admissible information and decision mechanism as their ReLU
encoding, so any difference must be attributed to actual implementation
work or separately proved information, not to this common inequality.

\paragraph{Finite computation account.}
Among $p$ paid records there are at most $\binom p2$ distinct ordered
brackets after omitting equal actions. Each contributes at most two
proposals and the neighbor rule at most six. Rounding and duplicate
removal cannot increase this count. A direct calculation of all
applicable pair uppers at $m$ candidates uses at most
$m\binom p2$ pair evaluations. The implementation evaluates the
origin bound again at each candidate; counting those evaluations
gives at most $2m\binom p2$ per completed stage. An uncompleted
stage retains all evaluations actually attempted. Evaluating one such expression and
its exact minimizer requires only fixed many rational operations;
their bit costs still depend on the stored rational lengths and are
not declared constant in a bit model. The original-risk calls,
fused calculations, failed/rejected attempts, optimization proposals,
exact dual checks, model acquisition and complete service clocks all
remain in the measured cost account. No deterministic policy-order
inequality implies an advantage after these computation charges.
